\documentclass[10pt,a4paper]{amsart}
\usepackage[margin=19mm]{geometry}
\usepackage{amsmath,amssymb,mathtools}
\usepackage[scr=rsfs,cal=euler]{mathalfa}
\usepackage{xfrac}
\usepackage[unicode,hidelinks,bookmarks=false]{hyperref}
\newcommand{\ie}{i.e.\ }
\newcommand{\cf}{cf.\ }
\newcommand{\resp}{respectively\ }
\newcommand{\Subsection}{\subsection}
\providecommand{\nobreakdash}{\nobreak\hbox{-}}
\setlength{\emergencystretch}{2em}
\multlinegap0pt
\usepackage{bm}
\usepackage{bbm}
\usepackage{dsfont}
\usepackage{xspace}
\usepackage{cancel}
\usepackage{mathtools}
\usepackage{cleveref}
\usepackage{amsthm}
\makeatletter
\@ifundefined{theorem}{\newtheorem{theorem}{Theorem}}{}
\@ifundefined{vmatrix}{\newtheorem{vmatrix}{Vmatrix}}{}
\makeatother
\title[Counting Frameworks of Bipyramids]{Counting Frameworks of Bipyramids}
\author{Jack Southgate}
\date{Source: arXiv:2301.04394v4 (CC BY 4.0)}
\begin{document}
\maketitle

\noindent\emph{Source and licence.} Counting Frameworks of Bipyramids. Version arXiv:2301.04394v4 (CC BY 4.0), distributed under the Creative Commons Attribution 4.0 International licence, as recorded in the arXiv licence metadata for this exact version and shown on the abstract page https://arxiv.org/abs/2301.04394v4. Full rights notice: licenses/arxiv-2301.04394-CC-BY-4.0.txt.\par\smallskip

\noindent\textit{Editorial scope.} Counted unit: the Theorem whose source label is \texttt{thm:bipcount} in the article (source0.tex lines 830--832, source label \texttt{thm:bipcount}), delivered together with the article's own complete original proof of it (source0.tex lines 834--951, one \texttt{proof} environment of 118 lines). No mathematics is written, repaired or re-derived: statement and proof are the article's own text line for line. Required context retained uncounted: none. Every definition, display and label of those lines is retained unchanged. Omitted from this package and not redistributed: 1--829, 833--833, 952--1069. Nothing inside a retained excerpt is omitted or altered: every retained body is delivered exactly as the article's lines are written, so no omission receipt of any kind accompanies this package. Citations: the retained excerpts carry no citation command at all, so no bibliography entry of the article is transcribed and no reference is redistributed. Reference adaptation: none was needed and none was made; every reference inside the delivered unit resolves inside it, so no unresolved reference or citation is produced. Printed slips kept literal: no printed word, symbol, hypothesis or number of the counted statement or of its proof was altered. Layout and class: the article's own class is \texttt{amsart} with packages including fontenc, lmodern, amsmath, amsthm, amssymb, amsfonts, enumerate, epsfig, xcolor, pgf, tikz, fullpage, hyperref, booktabs; the wrapper is the lane's pinned \texttt{amsart} editorial wrapper at 10pt on A4, which restates the theorem-like environments the retained text needs and adds the article's own macro definitions that the retained text needs (none), with extra wrapper packages (bm, bbm, dsfont, xspace, cancel, mathtools, cleveref). The delivered numbering is the unit's own. Assets: no retained span contains a figure environment, a graphics-inclusion command, a PDF-page inclusion, a table or a TikZ picture; no image file is copied into this package, the package declares no asset dependency and no image file is redistributed with it. Licence: version arXiv:2301.04394v4 of this article is distributed under the Creative Commons Attribution 4.0 International licence, as recorded in the arXiv licence metadata for this exact version and shown on the abstract page https://arxiv.org/abs/2301.04394v4. A full rights notice is delivered at licenses/arxiv-2301.04394-CC-BY-4.0.txt. Characters: every retained excerpt is pure ASCII, so the delivered engine, XeLaTeX with its default Latin Modern text font, reports no missing character.
\begin{theorem}\label{thm:bipcount}
    For each $n\geq 5$, $c(B_{n-2})\leq n-4$.
\end{theorem}

\begin{proof}
    To prove this theorem, we will take a generic framework $(B_{n-2},\mathbf{p})$, pin it to obtain $(B_{n-2},\overline{\mathbf{p}})$ and show that a pinned framework $(B_{n-2},\overline{\mathbf{q}})$ is equivalent to $(B_{n-2},\overline{\mathbf{p}})$ if and only if it satisfies a polynomial equation with coefficients in $\mathbb{Q}[\overline{\mathbf{p}}]$ of degree $n-4$.

    Pin $(B_{n-2},\mathbf{p})$ at 2-simplex $123$ so that $\overline{\mathbf{p}}(1)=(0,0)$, $\overline{\mathbf{p}}(2)=(1,0)$ and $\overline{\mathbf{p}}(3)=(0,1)$.
    Note again that, although this pinned framework is not generic, scaling by a generic quantity in the $y$-direction then performing a suitable special affine transformation of $\mathbb{R}^2$ yields a generic framework without changing any of its rigidity-theoretic properties.

    Suppose that $(B_{n-2},\overline{\mathbf{p}})$ is equivalent to $(B_{n-2},\overline{\mathbf{q}})$.
    Then, $\det(C(123,\overline{\mathbf{p}}))=\det(C(123,\overline{\mathbf{q}})$ (immediately, as they are both pinned to the same coordinates).
    The following three equalities define three lines on which the corresponding points of $(B_{n-2},\overline{\mathbf{q}})$ must lie:
    \begin{equation}
    \begin{split}
        \det(C(12(n-1),\overline{\mathbf{p}})) &= \det(C(12(n-1),\overline{\mathbf{p}})),\\
        \det(C(134,\overline{\mathbf{p}})) &= \det(C(134,\overline{\mathbf{p}})),\\
        \det(C(23n,\overline{\mathbf{p}})) &= \det(C(234,\overline{\mathbf{p}}));
    \end{split}
    \end{equation}
    respectively:
    \begin{equation}\label{eq:bip1}
        \overline{\mathbf{q}}(n-1) = \begin{bmatrix} \overline{\mathbf{p}}(n-1)_1 + s \\ \overline{\mathbf{p}}(n-1)_2 \end{bmatrix}\text{, }
        \overline{\mathbf{q}}(4) = \begin{bmatrix} \overline{\mathbf{p}}(4)_1 \\ \overline{\mathbf{p}}(4)_2 + t\end{bmatrix}\text{, }
        \overline{\mathbf{q}}(n) = \begin{bmatrix} \overline{\mathbf{p}}(n)_1 + r \\ \overline{\mathbf{p}}(n)_2 - r \end{bmatrix},
    \end{equation}
    with $r,s,t$ varying in $\mathbb{R}$.

    Next, plugging in the values from \cref{eq:bip1} into the equations $\det(C(2(n-1)n,\overline{\mathbf{p}}))=\det(C(2(n-1)n,\overline{\mathbf{q}}))$ and $\det(C(34n,\overline{\mathbf{p}}))=\det(C(34n,\overline{\mathbf{q}}))$, we obtain
    \begin{equation}\label{eq:bip2}
    \begin{split}
        r&=\frac{\overline{\mathbf{p}}(n)_1t}{1-\overline{\mathbf{p}}(4)_1-\overline{\mathbf{p}}(4)_2-t}=\frac{\overline{\mathbf{p}}(n)_2s}{1-\overline{\mathbf{p}}(n-1)_1-\overline{\mathbf{p}}(n-1)_2-2} \\
        \implies&s=\frac{(\overline{\mathbf{p}}(n-1)_1+\overline{\mathbf{p}}(n-1)_2-1)\overline{\mathbf{p}}(n)_1t}{(\overline{\mathbf{p}}(4)_1+\overline{\mathbf{p}}(4)_2-1)\overline{\mathbf{p}}(n)_2+(\overline{\mathbf{p}}(n)_1+\overline{\mathbf{p}}(n)_2)t},
    \end{split}
    \end{equation}
    thus putting all our equations in terms of one variable: $t$.

    Next, we define the positions of the equatorial vertices of $(B_{n-2},\overline{\mathbf{q}})$ in terms of $\overline{\mathbf{p}}$ and $t$.
    For each $4\leq i\leq n-1$, the position of $\overline{\mathbf{q}}(i)$, is completely determined by $\overline{\mathbf{q}}(1)=(0,0)$, $\overline{\mathbf{q}}(i-1)$ and $\overline{\mathbf{q}}(n)=(\overline{\mathbf{p}}(n)_1+r,\overline{\mathbf{p}}(n)_2-r)$.
    Therefore,
    \begin{equation}\label{eq:bip3}
    \begin{split}
        \overline{\mathbf{q}}(i)_j=&\frac{\left(\begin{vmatrix} \overline{\mathbf{q}}(i-1)_1 & \overline{\mathbf{q}}(n)_1 \\ \overline{\mathbf{q}}(i-1)_2 & \overline{\mathbf{q}}(n)_2 \end{vmatrix} - \begin{vmatrix} \overline{\mathbf{p}}(i-1)_1 & \overline{\mathbf{p}}(n)_1 \\ \overline{\mathbf{p}}(i-1)_2 & \overline{\mathbf{p}}(n)_2 \end{vmatrix} \begin{vmatrix} \overline{\mathbf{p}}(i)_1 & \overline{\mathbf{p}}(n)_1 \\ \overline{\mathbf{p}}(i)_2 & \overline{\mathbf{p}}(n)_2 \end{vmatrix} \right)\overline{\mathbf{q}}(i-1)_j}{\begin{vmatrix} \overline{\mathbf{q}}(i-1)_1 & \overline{\mathbf{q}}(n)_1 \\ \overline{\mathbf{q}}(i-1)_2 & \overline{\mathbf{q}}(n)_2 \end{vmatrix}} \\
        &+ \frac{\begin{vmatrix} \overline{\mathbf{p}}(i-1)_1 & \overline{\mathbf{p}}(n)_1 \\ \overline{\mathbf{p}}(i-1)_2 & \overline{\mathbf{p}}(n)_2 \end{vmatrix}\overline{\mathbf{q}}(n)_j}{\begin{vmatrix} \overline{\mathbf{q}}(i-1)_1 & \overline{\mathbf{q}}(n)_1 \\ \overline{\mathbf{q}}(i-1)_2 & \overline{\mathbf{q}}(n)_2 \end{vmatrix}},
    \end{split}
    \end{equation}
    for each $j\in[2]$.

    In order to simplify \cref{eq:bip3}, notice the following identities:
    \begin{equation}\label{eq:detind1}
        \begin{vmatrix} \overline{\mathbf{p}}(i)_1 & \overline{\mathbf{p}}(n)_1 \\ \overline{\mathbf{p}}(i-1)_2 & \overline{\mathbf{p}}(n)_2 \end{vmatrix} - \begin{vmatrix} \overline{\mathbf{p}}(i-1)_1 & \overline{\mathbf{p}}(n)_1 \\ \overline{\mathbf{p}}(i-1)_2 & \overline{\mathbf{p}}(n)_2 \end{vmatrix} = \det(C(1in,\overline{\mathbf{p}})) = \det(C(1(i-1)n,\overline{\mathbf{p}}))
    \end{equation}
    and
    \begin{equation}\label{eq:detind2}
        \det(C(1(i-1)i,\overline{\mathbf{p}})) - \det(C((i-1)in,\overline{\mathbf{p}})) = \det(C(1(i-1)n,\overline{\mathbf{p}})) - \det(C(1in,\overline{\mathbf{p}})).
    \end{equation}
    The right hand side of \cref{eq:detind2} is constant across equivalent pinned frameworks, hence
    \begin{equation}\label{eq:bip4}
        \begin{vmatrix} \overline{\mathbf{q}}(i)_1 & \overline{\mathbf{q}}(n)_1 \\ \overline{\mathbf{q}}(i)_2 & \overline{\mathbf{q}}(n)_2 \end{vmatrix} - \begin{vmatrix} \overline{\mathbf{q}}(i-1)_1 & \overline{\mathbf{q}}(n)_1 \\ \overline{\mathbf{q}}(i-1)_2 & \overline{\mathbf{q}}(n)_2 \end{vmatrix} = \begin{vmatrix} \overline{\mathbf{p}}(i)_1 & \overline{\mathbf{p}}(n)_1 \\ \overline{\mathbf{p}}(i)_2 & \overline{\mathbf{p}}(n)_2 \end{vmatrix} - \begin{vmatrix} \overline{\mathbf{p}}(i-1)_1 & \overline{\mathbf{p}}(n)_1 \\ \overline{\mathbf{p}}(i-1)_2 & \overline{\mathbf{p}}(n)_2 \end{vmatrix}
    \end{equation},
    for each $4\leq i\leq n-1$.
    Next, after applying \cref{eq:bip4} sufficiently many times, we obtain
    \begin{equation}
    \begin{split}
        \begin{vmatrix}
            \overline{\mathbf{q}}(i-1)_1 & \overline{\mathbf{q}}(n)_1 \\ \overline{\mathbf{q}}(i-1)_2 & \overline{\mathbf{q}}(n)_2
        \end{vmatrix} &= \begin{vmatrix} \overline{\mathbf{q}}(3)_1 & \overline{\mathbf{q}}(n)_1 \\ \overline{\mathbf{q}}(3)_2 & \overline{\mathbf{q}}(n)_2 \end{vmatrix} - \begin{vmatrix} \overline{\mathbf{p}}(3)_1 & \overline{\mathbf{p}}(n)_1 \\ \overline{\mathbf{p}}(3)_2 & \overline{\mathbf{p}}(n)_2 \end{vmatrix} + \begin{vmatrix} \overline{\mathbf{p}}(i-1)_1 & \overline{\mathbf{p}}(n)_1 \\ \overline{\mathbf{p}}(i-1)_2 & \overline{\mathbf{p}}(n)_2 \end{vmatrix} \\
        &= \begin{vmatrix} 0 & \overline{\mathbf{p}}(n)_1 + r \\ 1 & \overline{\mathbf{p}}(n)_2 - r \end{vmatrix} - \begin{vmatrix} 0 & \overline{\mathbf{p}}(n)_1 \\ 1 & \overline{\mathbf{p}}(n)_2 \end{vmatrix} + \begin{vmatrix} \overline{\mathbf{p}}(i-1)_1 & \overline{\mathbf{p}}(n)_1 \\ \overline{\mathbf{p}}(i-1)_2 & \overline{\mathbf{p}}(n)_2 \end{vmatrix} \\
        &= \begin{vmatrix}
            \overline{\mathbf{p}}(i-1)_1 & \overline{\mathbf{p}}(n)_1 \\ \overline{\mathbf{p}}(i-1)_2 & \overline{\mathbf{p}}(n)_2
        \end{vmatrix} - r.
    \end{split} 
    \end{equation}
    Plugging this into our formula from \cref{eq:bip3}, we obtain
    \begin{equation}\label{eq:bip5}
        \overline{\mathbf{q}}(i)_j = \frac{\left(\begin{vmatrix} \overline{\mathbf{p}}(i)_1 & \overline{\mathbf{p}}(n)_1 \\ \overline{\mathbf{p}}(i)_2 & \overline{\mathbf{p}}(n)_2 \end{vmatrix} - r\right)\overline{\mathbf{q}}(i-1)_j + \begin{vmatrix} \overline{\mathbf{p}}(i-1)_1 & \overline{\mathbf{p}}(i)_1 \\ \overline{\mathbf{p}}(i-1)_2 & \overline{\mathbf{p}}(i)_2 \end{vmatrix}(\overline{\mathbf{p}}(n)_j - (-1)^jr)}{\begin{vmatrix} \overline{\mathbf{p}}(i-1)_1 & \overline{\mathbf{p}}(i)_1 \\ \overline{\mathbf{p}}(i-1)_2 & \overline{\mathbf{p}}(i)_2 \end{vmatrix} - r},
    \end{equation}
    for each $j\in[2]$ and $4\leq i \leq n-1$.

    In order for $(B_{n-2},\overline{\mathbf{q}})$ as we are defining it to be a well-defined framework, the position of $\overline{\mathbf{q}}(n-1)$ obtained from \cref{eq:bip5} must line up with that from \cref{eq:bip1}.
    This equality is obtained when the following two equations are satisfied:
    \begin{equation}\label{eq:bip6}
    \begin{split}
        &\left(\left( \begin{vmatrix} \overline{\mathbf{p}}(n-1)_1 & \overline{\mathbf{p}}(n)_1 \\ \overline{\mathbf{p}}(n-1)_2 & \overline{\mathbf{p}}(n)_2 \end{vmatrix} (1-\overline{\mathbf{p}}(4)_1 - \overline{\mathbf{p}}(4)_2 -t) - \overline{\mathbf{p}}(n)_1t \right)\overline{\mathbf{q}}(n-2)_1\right.\\
        &+\left. \begin{vmatrix} \overline{\mathbf{p}}(n-2)_1 & \overline{\mathbf{p}}(n-1)_1 \\ \overline{\mathbf{p}}(n-2)_2 & \overline{\mathbf{p}}(n-2)_2 \end{vmatrix} (1-\overline{\mathbf{p}}(4)_1-\overline{\mathbf{p}}(4)_2)\overline{\mathbf{p}}(n)_1 \right) \\
        &((\overline{\mathbf{p}}(4)_1+\overline{\mathbf{p}}(4)_2 - 1)\overline{\mathbf{p}}(n)_2 - (\overline{\mathbf{p}}(n)_1+\overline{\mathbf{p}}(n)_2)t) \\
        =& \left( \begin{vmatrix} \overline{\mathbf{p}}(n-2)_1 & \overline{\mathbf{p}}(n-1)_1 \\ \overline{\mathbf{p}}(n-2)_2 & \overline{\mathbf{p}}(n-1)_2 \end{vmatrix} (1-\overline{\mathbf{p}}(4)_1-\overline{\mathbf{p}}(4)_2-t) - \overline{\mathbf{p}}(n)_1t \right) \\
        &\left( (\overline{\mathbf{p}}(4)_1+\overline{\mathbf{p}}(4)_2-1)\overline{\mathbf{p}}(n-1)_1\overline{\mathbf{p}}(n)_2 + \left(\begin{vmatrix} \overline{\mathbf{p}}(n-1)_1 & \overline{\mathbf{p}}(n)_1 \\ \overline{\mathbf{p}}(n-1)_2 & \overline{\mathbf{p}}(n)_2 \end{vmatrix}\right)t \right)
    \end{split}
    \end{equation}
    and
    \begin{equation}\label{eq:bip7}
    \begin{split}
        &\left(\left( \begin{vmatrix} \overline{\mathbf{p}}(n-1)_1 & \overline{\mathbf{p}}(n)_1 \\ \overline{\mathbf{p}}(n-1)_2 & \overline{\mathbf{p}}(n)_2 \end{vmatrix} (1-\overline{\mathbf{p}}(4)_1 - \overline{\mathbf{p}}(4)_2 -t) - \overline{\mathbf{p}}(n)_1t \right)\overline{\mathbf{q}}(n-2)_2\right.\\
        &+\left. \begin{vmatrix} \overline{\mathbf{p}}(n-2)_1 & \overline{\mathbf{p}}(n-1)_1 \\ \overline{\mathbf{p}}(n-2)_2 & \overline{\mathbf{p}}(n-2)_2 \end{vmatrix} ((1-\overline{\mathbf{p}}(4)_1-\overline{\mathbf{p}}(4)_2)\overline{\mathbf{p}}(n)_2 - \overline{\mathbf{p}}(n)_1t) \right) \\
        =& \left( \begin{vmatrix} \overline{\mathbf{p}}(n-2)_1 & \overline{\mathbf{p}}(n-1)_1 \\ \overline{\mathbf{p}}(n-2)_2 & \overline{\mathbf{p}}(n-1)_2 \end{vmatrix} (1-\overline{\mathbf{p}}(4)_1-\overline{\mathbf{p}}(4)_2-t) - \overline{\mathbf{p}}(n)_1t \right) \overline{\mathbf{p}}(n-1)_2.
    \end{split}
    \end{equation}
    Notice, however, that it suffices to solve just \cref{eq:bip7}.
    Indeed, if $\tau$ solves \cref{eq:bip7}, so $\overline{\mathbf{q}}(\tau)(n-1)_2=\overline{\mathbf{p}}(n-1)_2$, but not \cref{eq:bip6}, so $\overline{\mathbf{q}}(\tau)(n-1)_1\neq\overline{\mathbf{p}}(n-1)_1+\sigma$, for some $\sigma\neq s$, then comparing the value of $r$ yielded by $\det(C(2(n-1)n,\overline{\mathbf{p}}))=\det(C(2(n-1)n,\overline{\mathbf{q}}(\tau)))$ gives us
    \begin{equation*}
        s(\overline{\mathbf{p}}(n-1)_1+\overline{\mathbf{p}}(n-1)_2-1)=\sigma(\overline{\mathbf{p}}(n-1)_1+\overline{\mathbf{p}}(n-1)_2),
    \end{equation*}
    hence $s=\sigma$, a contradiction.

    Subtracting the right hand side from both sides of \cref{eq:bip7} yields an equation of the form $\text{LHS}=0$, then multiplying through by the denominator of $\overline{\mathbf{q}}(n-2)_2$ increases the degree of terms not containing $\overline{\mathbf{q}}(n-2)_2$ by 1 (with respect to the variable $t$).
    This yields an equation of the form $f(\overline{\mathbf{p}})(t)=0$, with
    \begin{equation*}
        \deg(f(\mathbf{p}))=\begin{cases}
            \max\{\deg(\overline{\mathbf{q}}(n-2)_2)+1,2\},\text{ if }n\geq 6,\\
            \deg(\overline{\mathbf{q}}(n-2)_2+1,\text{ if }n=5.
        \end{cases}
    \end{equation*}
    By considering \cref{eq:bip5}, we notice that the degree of the numerator of $\overline{\mathbf{q}}(i)_2$ has a higher degree than its denominator, and so, multiplying $f(\mathbf{p})(t)=0$ through by the denominator of $\overline{\mathbf{q}}(n-2)_2$ does not change the degree of $f(\mathbf{p})$, hence the numerator's degree increases by 1 for each equatorial vertex, with an initial value of 0 at $\overline{\mathbf{q}}(3)_2$.
    Therefore $\deg(f(\mathbf{p}))=n-4$.

    This completes the proof subject to one final claims:
        If $(B_{n-2},\mathbf{p})$ and $(B_{n-2},\mathbf{q}_1)$ and $(B_{n-2},\mathbf{q}_2)$ are equivalent and if $(B_{n-2},\overline{\mathbf{q}_1})\ne\mathbf{q}((B_{n-2},\overline{\mathbf{q}_2})$, then $t_1=\overline{\mathbf{q}_1}(4)_2-\overline{\mathbf{p}}(4)_2\neq\overline{\mathbf{q}_2}(4)_2-\overline{\mathbf{p}}(4)_2=t_2$.

        Indeed, if, to the contrary, $t_1=t_2$, then $r_1=r_2$ and $s_1=s_2$ (with $r_1,r_2,s_1,s_2$ analogously defined).
        Therefore each equatorial vertex of $(B_{n-2},\overline{\mathbf{q}_2})$ has the same position of its corresponding vertex in $(B_{n-2},\overline{\mathbf{q}_1})$, hence $(B_{n-2},\overline{\mathbf{q}_1})=(B_{n-2},\overline{\mathbf{q}_2})$.
\end{proof}

\end{document}
